% Problems numbered within each section. The setting holds until it is
% changed again, so it is placed before the chapter; chapters 1 and 2 keep
% the default numbering within chapters.
\SolutionManualSetup{problem numbering=subdivision}

\chapter{Sections and Progress}

% All the problems of a chapter are at the same level: here, in sections.

\section{Sequences}

% ============ Section 3.1, Problem 1 ============ %
% A completed problem. Its references: 3.1.1 and \probref{3.1}{1}.

\begin{problem}
{
Compute the sum $1 + 2 + 3 + \dots + 100$.
}
{
Pair the first term with the last, the second with the second to last, and
so on: each of the $50$ pairs $k + (101 - k)$ adds up to $101$.
}
{
$1 + 2 + \dots + 100 = 50 \cdot 101 = 5050$.
}
\end{problem}

% ============ Section 3.1, Problem 2 ============ %

\begin{problem}
{
An arithmetic sequence starts with $3, 7, 11, \dots$. Find its tenth term
$a_{10}$ and the sum of its first ten terms.\label{prob:arithmetic}
}
{
The common difference is $4$, so $a_n = 3 + 4(n - 1)$ and $a_{10} = 39$. The
sum of an arithmetic sequence is the number of terms times the mean of the
first and last terms:
\begin{align}
    a_1 + \dots + a_n &= \frac{n (a_1 + a_n)}{2}. \label{eq:arithmetic-sum}
\end{align}
}
{
$a_{10} = 39$, and the sum is $\frac{10 (3 + 39)}{2} = 210$.
}
\end{problem}

% ============ Section 3.1, Problem 3 ============ %
% An unanswered question: no solution, so it is not done.

\begin{problem}
{
Find a closed form for the sum of the first $n$ squares,
$1^2 + 2^2 + \dots + n^2$.
}{}{}
\end{problem}

% ============ Section 3.1, Problem 4 ============ %
% A partial solution, explicitly marked [not done].

\begin{problem}[not done]
{
Show that the sequence $a_n = \left( 1 + \frac{1}{n} \right)^n$ is
increasing and bounded above by $3$.
}
{
By the binomial theorem,
\begin{align*}
    a_n = \sum_{k=0}^{n} \binom{n}{k} \frac{1}{n^k}
        = \sum_{k=0}^{n} \frac{1}{k!} \prod_{j=0}^{k-1} \left( 1 - \frac{j}{n} \right).
\end{align*}
Each product grows with $n$, and $a_{n+1}$ has one more positive term than
$a_n$, so the sequence is increasing. The bound remains to be proved.
}
{}
\end{problem}

% A section shown with the number of the textbook (§12) instead of the
% automatic one (3.2). Its problems restart at 1: references §12.1, ...
\textbooknumber{\S 12}
\section{Series}

% ============ §12, Problem 1 ============ %

\begin{problem}
{
Find the sum of the geometric series $1 + \frac{1}{2} + \frac{1}{4} + \dots$.
\label{prob:geometric}
}
{
The partial sums are
\begin{align*}
    S_n = \sum_{k=0}^{n-1} \frac{1}{2^k}
        = \frac{1 - 2^{-n}}{1 - \frac{1}{2}}
        = 2 \left( 1 - 2^{-n} \right),
\end{align*}
and $2^{-n} \to 0$.
}
{
The sum is $2$.
}
\end{problem}

% ============ §12, Problem 2 ============ %
% A problem in parts, done only when all its parts are done: (a) is done,
% (b) is not (no solution), and (c) gives only an answer, marked [done].

\begin{problem}
{
Consider the series $\displaystyle \sum_{k=1}^{\infty} \frac{1}{k (k + 1)}$.
\label{prob:telescoping}
}{}{}
\end{problem}

\begin{subproblem}
{
Write $\frac{1}{k (k + 1)}$ as a difference of two simple fractions.
\label{prob:telescoping-a}
}
{
Since $(k + 1) - k = 1$,
\begin{align*}
    \frac{1}{k} - \frac{1}{k + 1} = \frac{(k + 1) - k}{k (k + 1)}
                                  = \frac{1}{k (k + 1)}.
\end{align*}
}
{
$\frac{1}{k (k + 1)} = \frac{1}{k} - \frac{1}{k + 1}$.
}
\end{subproblem}

\begin{subproblem}
{
Compute the partial sum $S_n = \sum_{k=1}^{n} \frac{1}{k (k + 1)}$.
}{}{}
\end{subproblem}

\begin{subproblem}[done]
{
Deduce the sum of the series.
}
{}
{
The sum is $1$.
}
\end{subproblem}

% ============ §12, Problem 3 ============ %

\begin{problem}
{
Show that the series $\displaystyle \sum_{k=1}^{\infty} \frac{1}{k^2}$
converges.
}
{
For $k \ge 2$, $\frac{1}{k^2} < \frac{1}{(k - 1) k}$. By part
\ref{prob:telescoping-a} (\probref{\S 12}{2}), these terms
telescope:
\begin{align*}
    \sum_{k=1}^{n} \frac{1}{k^2}
        < 1 + \sum_{k=2}^{n} \left( \frac{1}{k - 1} - \frac{1}{k} \right)
        = 2 - \frac{1}{n} < 2.
\end{align*}
The partial sums increase and are bounded, so the series converges.
}
{
The series converges, and its sum is less than $2$.
}
\end{problem}

% ============ §12, Problem 4 ============ %
% References: \ref gives the full number (3.1.2, §12.1); \probref links to a
% problem from the number of its chapter or section and its own number.

\begin{problem}
{
Find the sum of the first $20$ terms of the sequence of
Problem~\ref{prob:arithmetic}, and compare it with the sum of the first $20$
odd numbers.
}
{
By \eqref{eq:arithmetic-sum} (\probref{3.1}{2}), with $a_{20} = 3 + 4
\cdot 19 = 79$, the sum is $\frac{20 (3 + 79)}{2} = 820$. The sum of the first
$20$ odd numbers is $20^2 = 400$ (\probref{1}{2}), and the series of
Problem~\ref{prob:geometric} shows that sums of infinitely many terms can
still be finite.
}
{
$820$, against $400$ for the odd numbers.
}
\end{problem}
